\section{Documentation technique}

\subsection{Architecture générale}

Le code est organisé autour d'une boucle principale située dans
\texttt{src/systus/app.py}. Cette boucle initialise le mode système, démarre le
serveur Wi-Fi local, configure les boutons, puis met à jour le contrôleur de
detection.

\begin{lstlisting}
src/systus/
  app.py
  audio_capture/
    recorder.py
  buttons/
    controller.py
  detection/
    audd_client.py
    detection.py
  display/
    screen.py
  wifi/
    wifi.py
    templates/
    static/
\end{lstlisting}

\subsection{Modes système}

Deux modes structurent le comportement :

\begin{longtable}{|p{3cm}|p{10cm}|}
\hline
\textbf{Mode} & \textbf{Rôle} \\
\hline
\texttt{setup} & Mode de configuration. L'écran affiche un QR code permettant d'accéder à l'interface Wi-Fi locale. \\
\hline
\texttt{running} & Mode d'utilisation normale. L'utilisateur peut lancer une reconnaissance musicale. \\
\hline
\end{longtable}

Au démarrage, le programme teste l'état réseau avec \texttt{nmcli}. Si le Wi-Fi
est connecté, le mode \texttt{running} est choisi ; sinon, le mode
\texttt{setup} est affiché.

\subsection{Capture audio}

La capture audio est assurée par la classe \texttt{AudioRecorder}. Elle crée un
fichier temporaire au format WAV, lance la commande \texttt{arecord}, puis
retourne le chemin du fichier capturé.

Paramètres utilisés :

\begin{itemize}
    \item périphérique : \texttt{plughw:0,0} ;
    \item format : \texttt{S16\_LE} ;
    \item fréquence : 16000 Hz ;
    \item canal : mono ;
    \item durée de détection actuelle : 10 secondes.
\end{itemize}

Extrait logique :

\begin{lstlisting}
arecord -D plughw:0,0 -f S16_LE -r 16000 -c 1 -d 10 fichier.wav
\end{lstlisting}

\subsection{Reconnaissance musicale}

La reconnaissance est assurée par la classe \texttt{AudDClient}. Elle envoie le
fichier WAV à l'API AudD avec une requête HTTP \texttt{POST}. La réponse est
ensuite vérifiée : si le statut est \texttt{success}, le champ \texttt{result}
est conservé.

Métadonnées exploitées par l'interface :

\begin{itemize}
    \item titre ;
    \item artiste ;
    \item album ;
    \item date de sortie ;
    \item lien vers le morceau.
\end{itemize}

Limite importante : dans l'état actuel du code, le token AudD est écrit en dur
dans le contrôleur de détection. Pour une version propre, il doit être déplacé
dans une variable d'environnement ou un fichier de configuration non commité.

\subsection{Contrôleur de détection}

Le module \texttt{detection.py} contient une machine à états simple :

\begin{longtable}{|p{3cm}|p{10cm}|}
\hline
\textbf{État} & \textbf{Description} \\
\hline
\texttt{idle} & Le système attend une action utilisateur. \\
\hline
\texttt{listening} & Le système enregistre un extrait audio. \\
\hline
\texttt{thinking} & Le système envoie l'audio à AudD et attend le résultat. \\
\hline
\texttt{result} & Le résultat ou l'échec est affiché. \\
\hline
\end{longtable}

Cette structure évite de mélanger la logique de capture, la reconnaissance et
l'affichage. Elle rend aussi les états faciles à expliquer lors de la
soutenance.

\subsection{Affichage e-paper}

Le module \texttt{screen.py} utilise la bibliothèque e-paper Waveshare, ainsi
que \texttt{Pillow} et \texttt{qrcode}. L'écran est traité comme une image
monochrome de 400 par 300 pixels.

Vues disponibles :

\begin{itemize}
    \item \texttt{show\_setup} : mode configuration avec QR code ;
    \item \texttt{show\_idle} : attente d'une détection ;
    \item \texttt{show\_listening} : capture en cours ;
    \item \texttt{show\_thinking} : reconnaissance en cours ;
    \item \texttt{show\_result} : affichage titre, artiste, album, année et QR code.
\end{itemize}

\subsection{Interface Wi-Fi}

Le module \texttt{wifi.py} crée une application Flask. Elle expose trois routes
principales :

\begin{longtable}{|p{3cm}|p{3cm}|p{7cm}|}
\hline
\textbf{Route} & \textbf{Méthode} & \textbf{Rôle} \\
\hline
\texttt{/} & GET & Affiche la page de configuration avec les réseaux détectés. \\
\hline
\texttt{/scan} & GET & Retourne les réseaux Wi-Fi disponibles en JSON. \\
\hline
\texttt{/wifi} & POST & Ajoute/modifie une connexion Wi-Fi via \texttt{nmcli}. \\
\hline
\end{longtable}

La configuration réseau repose sur \texttt{nmcli}. Le système supprime
d'abord une éventuelle ancienne connexion du même nom, recrée la connexion,
configure la sécurité WPA-PSK, puis tente de l'activer.

\subsection{Boutons physiques}

Le module \texttt{controller.py} utilise \texttt{gpiozero}. Deux boutons sont
déclarés :

\begin{itemize}
    \item GPIO 5 : changement de mode, notamment accès au mode de mise en place du Wi-Fi ;
    \item GPIO 6 : action principale, lancement manuel d'une écoute.
\end{itemize}

Un verrou et une temporisation de 0,4 seconde limitent les doubles déclenchements
sur le bouton de mode. L'écoute n'est pas continue : chaque tentative de
reconnaissance doit être lancée par un nouvel appui sur le bouton d'action.

\subsection{Matériel identifié}

\begin{longtable}{|p{5cm}|p{4cm}|p{4cm}|}
\hline
\textbf{Élément} & \textbf{Statut} & \textbf{Remarque} \\
\hline
Raspberry Pi Zero 2 W & Intégré au choix matériel & Version avec headers soudés. \\
\hline
Pico-ePaper-4.2, 400 x 300 px & Prévu / codé & Écran noir/blanc avec 4 niveaux de gris. \\
\hline
Microphone MEMS ICS-43434 I2S & Intégré au choix matériel & Module 3,3 V. \\
\hline
Boutons physiques & Prévu / codé & GPIO 5 et GPIO 6. \\
\hline
Batterie LiPo 2000 mAh, 7,4 Wh & Intégrée au choix matériel & Connecteur JST PH 2.0 mm, 2 broches. \\
\hline
UPS HAT (C) pour Raspberry Pi Zero & Intégré au choix matériel & Version sans batterie fournie. \\
\hline
Boîtier / support / pièce écran & Archivé & Fichiers STL dans \texttt{docs/assets/cad}. \\
\hline
\end{longtable}

\subsection{Schéma électrique}

Le schéma électrique fourni relie le Raspberry Pi Zero 2 W au microphone
ICS-43434, à l'écran e-paper, aux boutons et à l'alimentation. Il fait apparaître
les rails 5 V et 3,3 V, le connecteur du microphone I2S, le connecteur de
l'écran e-paper et trois interrupteurs momentanés prévus pour les interactions
physiques.

Dans le code actuel, deux entrées bouton sont utilisées :

\begin{itemize}
    \item GPIO 5 pour la bascule de mode ;
    \item GPIO 6 pour le lancement d'écoute.
\end{itemize}

Le troisième bouton visible sur le schéma peut être conservé pour une évolution
future, par exemple changement d'écran, mise en veille ou relance rapide.

\subsection{Budget prévisionnel}

\begin{longtable}{|p{5cm}|p{2cm}|p{3cm}|p{3cm}|}
\hline
\textbf{Composant} & \textbf{Qté} & \textbf{Prix unitaire} & \textbf{Total} \\
\hline
Raspberry Pi Zero 2 W avec headers & 1 & 54,99 \texteuro{} & 54,99 \texteuro{} \\
\hline
Pico-ePaper-4.2, 400 x 300 & 1 & 28,19 \texteuro{} & 28,19 \texteuro{} \\
\hline
Microphone MEMS ICS-43434 I2S & 1 & 7,39 \texteuro{} & 7,39 \texteuro{} \\
\hline
Lot de 10 boutons poussoirs PB86-A0 & 1 & 6,49 \texteuro{} & 6,49 \texteuro{} \\
\hline
Batterie LiPo 3,7 V 2000 mAh & 1 & 6,79 \texteuro{} & 6,79 \texteuro{} \\
\hline
UPS HAT (C) Raspberry Pi Zero & 1 & 24,59 \texteuro{} & 24,59 \texteuro{} \\
\hline
Carte microSD SanDisk 16 Go & 1 & 20,39 \texteuro{} & 20,39 \texteuro{} \\
\hline
\textbf{Total composants commandés} & & & \textbf{148,83 \texteuro{}} \\
\hline
\end{longtable}

Les commandes listées ont été passées le 28 avril 2026. Les frais éventuels de
matière pour l'impression 3D, câbles, visserie ou consommables ne sont pas
inclus dans ce total.
